% LaTeX document detailing sloshing FEM theory
\documentclass{article}
\usepackage{amsmath, amssymb, bm, geometry}
\geometry{margin=2cm}
\usepackage{tikz}

\title{Finite Element Formulation for Linear Sloshing in a Rectangular Tank}
\author{}
\date{}

\begin{document}
\maketitle

\section{Introduction}
This document outlines the theoretical foundations for the finite element method (FEM) applied to linear sloshing of an incompressible, inviscid fluid inside a rigid container. The formulation is based on potential flow theory and its linearization for small-amplitude free-surface oscillations.

\section{Governing Equations}
For an incompressible and inviscid fluid, the velocity field is expressed through a scalar potential:
\begin{equation}
    \bm{v} = \nabla \phi,
\end{equation}
where $\phi(x,y,t)$ is the velocity potential. Incompressibility requires
\begin{equation}
    \nabla^2 \phi = 0 \quad \text{in the fluid domain}.
\end{equation}

The pressure field follows from the linearized Bernoulli equation:
\begin{equation}
    p' = -\rho \frac{\partial \phi}{\partial t} - \rho g \, \eta,
\end{equation}
where $\eta(x,t)$ is the vertical displacement of the free surface.

\section{Boundary Conditions}
\subsection{Rigid Walls}
On the walls and bottom, the fluid cannot cross the boundary:
\begin{equation}
    \frac{\partial \phi}{\partial n} = 0.
\end{equation}

\subsection{Linearized Free-Surface Boundary Conditions}
Two boundary conditions apply on the free surface ($y=H$):

\begin{enumerate}
    \item Kinematic condition:
    \begin{equation}
        \frac{\partial \eta}{\partial t} = \frac{\partial \phi}{\partial y}.
    \end{equation}

    \item Dynamic condition:
    \begin{equation}
        \frac{\partial \phi}{\partial t} + g \, \eta = 0.
    \end{equation}
\end{enumerate}

Combining the two yields a single evolution equation for $\eta$:
\begin{equation}
    \frac{\partial^2 \eta}{\partial t^2} + g \, \frac{\partial \phi}{\partial y} = 0.
\end{equation}

\section{Reduction to a Free-Surface FEM Problem}
For small-amplitude sloshing, it is common to eliminate $\phi$ and derive a boundary integral or a reduced free-surface FEM equation. A classical approach leads to the eigenvalue problem:

\begin{equation}
    \bm{K} \bm{\eta} = \omega^2 \bm{M} \bm{\eta},
\end{equation}
where $\bm{\eta}$ is the nodal vector of free-surface elevations.

The matrix $\bm{K}$ comes from gravitational restoring effects:
\begin{equation}
    K_{ij} = \rho g \int_{A_e} \nabla N_i \cdot \nabla N_j ~ dA,
\end{equation}
with $N_i$ the shape functions of triangular CST elements.

The mass matrix, assuming lumping, is
\begin{equation}
    M_{ii} = \rho \frac{A_e}{3},
\end{equation}
for each node of a triangular element of area $A_e$.

\section{Element Matrices}
For a triangular element with nodes $(i,j,k)$, area
\begin{equation}
    A = \frac{1}{2} | x_i(y_j - y_k) + x_j(y_k - y_i) + x_k(y_i - y_j) |,
\end{equation}
the gradients of the shape functions are expressed using standard CST coefficients:
\begin{align}
    b_1 &= y_j - y_k, & c_1 &= x_k - x_j, \\
    b_2 &= y_k - y_i, & c_2 &= x_i - x_k, \\
    b_3 &= y_i - y_j, & c_3 &= x_j - x_i.
\end{align}

The stiffness matrix becomes:
\begin{equation}
    K_{ab} = \rho g \, t \frac{b_a b_b + c_a c_b}{4 A},
\end{equation}
where $t$ is the out-of-plane thickness.

The lumped mass matrix is:
\begin{equation}
    M = \rho t A \begin{bmatrix}
        1/3 & 0 & 0 \\
        0 & 1/3 & 0 \\
        0 & 0 & 1/3
    \end{bmatrix}.
\end{equation}

\section{Eigenvalue Problem and Modal Analysis}
The generalized eigenproblem
\begin{equation}
    \bm{K} \bm{\eta} = \omega^2 \bm{M} \bm{\eta}
\end{equation}
leads to natural frequencies
\begin{equation}
    f = \frac{\omega}{2 \pi}
\end{equation}
and periods
\begin{equation}
    T = \frac{1}{f}.
\end{equation}

The eigenvectors describe the shape of the free-surface oscillations. The pressure fluctuations inside the fluid follow:
\begin{equation}
    p'(x,y) = \rho g \, \eta(x).
\end{equation}

\section{References}
\begin{itemize}
    \item Faltinsen, O. M., \textit{Sea Loads on Ships and Offshore Structures}. Cambridge Univ. Press, 1990.
    \item Mei, C. C., \textit{The Applied Dynamics of Ocean Surface Waves}. World Scientific, 1989.
    \item Bathe, K.-J., \textit{Finite Element Procedures}. Prentice Hall, 1996.
    \item Zienkiewicz, O. C., Taylor, R. L., \textit{The Finite Element Method}. McGraw-Hill.
\end{itemize}

\textcolor{red}{\Huge VERSIONE SUCCESSIVA}

\section{Governing Equations (Expanded)}
For an incompressible, inviscid fluid, the flow is irrotational and can be described by a velocity potential $\phi(x,y,t)$:
\begin{equation}
\bm{v} = \nabla \phi,
\end{equation}
where $\bm{v}$ is the fluid velocity vector. The incompressibility condition leads to Laplace's equation in the fluid domain $\Omega$:
\begin{equation}
\nabla^2 \phi = \frac{\partial^2 \phi}{\partial x^2} + \frac{\partial^2 \phi}{\partial y^2} = 0, \quad \text{in } \Omega.
\end{equation}


On the rigid boundaries (walls and bottom), the normal component of velocity must vanish:
\begin{equation}
\frac{\partial \phi}{\partial n} = 0.
\end{equation}


At the free surface $y = H$, the boundary conditions are linearized for small-amplitude oscillations:
\begin{enumerate}
\item Kinematic condition:
\begin{equation}
\frac{\partial \eta}{\partial t} = \frac{\partial \phi}{\partial y} \Big|_{y=H},
\end{equation}
where $\eta(x,t)$ is the vertical displacement of the free surface.


\item Dynamic condition:
\begin{equation}
\frac{\partial \phi}{\partial t} + g \, \eta = 0, \quad y = H.
\end{equation}
\end{enumerate}


By differentiating the dynamic condition with respect to time and substituting the kinematic condition, one obtains a single second-order differential equation for the free-surface displacement $\eta(x,t)$:
\begin{equation}
\frac{\partial^2 \eta}{\partial t^2} + g \, \frac{\partial \phi}{\partial y} \Big|_{y=H} = 0.
\end{equation}


In a variational (FEM) formulation, this is typically converted into a weak form over the fluid domain, leading to a generalized eigenvalue problem for the nodal free-surface elevations:
\begin{equation}
\bm{K} \bm{\eta} = \omega^2 \bm{M} \bm{\eta},
\end{equation}
where $\bm{\eta}$ is the vector of nodal elevations (spostamento verticale del pelo libero), $\bm{M}$ is the mass matrix, and $\bm{K}$ is the stiffness matrix derived from the gravitational restoring force.


\section{Connection Between PDE and FEM Matrices}
The stiffness matrix $\bm{K}$ represents the discrete version of the linear operator associated with the free-surface restoring force:
\begin{equation}
K_{ij} = \rho g \int_{A_e} \nabla N_i \cdot \nabla N_j ~ dA,
\end{equation}
where $N_i$ and $N_j$ are the element shape functions, and $A_e$ is the element area. Essentially, $\nabla N_i \cdot \nabla N_j$ arises from the weak form of the Laplace equation projected onto the free-surface displacement field $\eta$.


The mass matrix $\bm{M}$ represents the inertial term $\partial^2 \eta/\partial t^2$ in the PDE, typically assembled in lumped or consistent form:
\begin{equation}
M_{ii} = \rho t A_e / 3 \quad \text{for each node of a triangular element}.
\end{equation}


Thus, the FEM eigenvalue problem directly approximates the PDE for the free-surface displacement. The eigenvectors provide the mode shapes of the surface, and the eigenvalues determine the natural sloshing frequencies.









\end{document}